% Source: https://artofproblemsolving.com/community/c3h1598561_2011_mathcounts_national_round
% Stable identity: aops-topic-1598561-q06
Cross-multiplying (both denominators are nonzero: $a\ne0$ and $a\ne b$),
\[
(a-b)^2=ab\qquad\Longleftrightarrow\qquad a^2-3ab+b^2=0 .
\]
Dividing by $b^2\ne0$ and writing $r=\dfrac {a}{b}$:
\[
r^2-3r+1=0 .
\]
Its discriminant $9-4=5$ is positive, so there are two real roots $r=\dfrac{3\pm\sqrt5}{2}$. Both are allowed: neither is $0$ or $1$, so $a\ne0$ and $a\ne b$ for every $b\ne0$, and the original equation holds. By Vieta's formulas the sum of the possible values is $\boxed{3}$.
